\section{Models retaining finitely many scalar tests}\label{sec:models}

Let $B\subset\C$ be compact and totally disconnected, and put
$G=\Chat\setminus B$.  For a normalized univalent map $h$ on $G$,
Lemma~\ref{lem:ends} gives a continuous collapse map
$\pi_h:\Chat\to\Chat$: it equals $h^{-1}$ on $h(G)$ and sends the
complementary component corresponding to $b\in B$ to $b$.

\begin{definition}\label{def:model}
Let $u_1,\ldots,u_m$ be bounded continuous real functions on $\Chat$
whose restrictions to $G$ belong to $W^{1,2}(G)$.  A \emph{model} for
this list is a univalent conformal map
\[
 h(z)=z+\frac{b(h)}z+O(z^{-2})\qquad(z\to\infty)
\]
such that its nonpoint complementary components form a finite or
countable null family of closed Jordan disks, and
\[
 U_i=u_i\circ\pi_h\in W^{1,2}(\Chat),\qquad
 \nabla U_i=0\quad\text{a.e. on }\Chat\setminus h(G).
\]
A family is null if only finitely many members have diameter greater
than any prescribed positive number.
\end{definition}

The conductivity argument below postcomposes models with quasiconformal
homeomorphisms.  These preserve Jordan disks and null families but need
not preserve circles, which is why the model class allows Jordan disks.
The existence result below will nevertheless supply round disks.

We require the pairwise trace condition
\begin{equation}\label{eq:pair-null}
 \area\bigl((u_i,u_j)(B)\bigr)=0\qquad(i\ne j).
\end{equation}
Only the values on $B$ occur in this condition.

\begin{proposition}\label{prop:models}
Every finite list in Definition~\ref{def:model} satisfying
\eqref{eq:pair-null} has a model.  A model can, in fact, be chosen with
circle-domain image.
\end{proposition}

The uniformization input below is the finite-transfer construction of
\cite{KoebeCompanion}.  Its control of the retained scalar tests is
essential; circle-domain existence by itself would not give
Proposition~\ref{prop:models}.

\subsection{The precise retained-test input}

We use the following specialization of
\cite[\CompanionRef{uniformization}{cor:scalar-retained-tests}]{KoebeCompanion}.
It is a consequence of the companion's
\CompanionRef{uniformization}{thm:retained-tests}, whose selection begins
with the prescribed tests and preserves them through all later compatible
extensions.

First suppose the tests $u_1,\ldots,u_m$ in
Proposition~\ref{prop:models} are smooth on $G$.  For any sequence
$\varepsilon_n>0$ tending to zero, there are nested compact smooth cores
$K_n\subset G$, finite bordered completions $Q_n$ containing
neighborhoods of $K_n$ with their original conformal structure, normalized conformal
maps $f_n:Q_n^\circ\to\Omega_n$ onto finite circle domains, and smooth
transferred tests $v_{i,n}$ on $\Omega_n$ with the following properties.
Every compact subset of $G$ is contained in $\operatorname{int}K_n$ for
all sufficiently large $n$, and
\begin{align}
 v_{i,n}\circ f_n&=u_i
       &&\text{near }K_n,\label{eq:model-core}\\
 \int_{\Omega_n\setminus f_n(K_n)}|\nabla v_{i,n}|^2
       &\le\varepsilon_n^2,\label{eq:model-tail}\\
 \sum_T A_{i,n}(T)^2&\le\varepsilon_n^2.
       \label{eq:model-osc}
\end{align}
Here $T$ ranges over the complementary round disks of $\Omega_n$, and
the cluster values of $v_{i,n}$ at $T$ lie in an interval of length
$A_{i,n}(T)$.  The first equality also preserves energy near $K_n$.  The
companion actually bounds the sum of the tail energies over the whole
retained finite list by $\varepsilon_n^2$.

These assertions hold on one shared sequence.  In particular, every
fixed source contour and its collar are retained eventually.  If $C$
is an oriented source Jordan contour and $z\in K_n$ lies on its
bounded side, then $f_n(z)$ lies on the bounded side of $f_n(C)$;
the converse holds for retained-core points.  After passage to a
subsequence,
\begin{equation}\label{eq:model-map-limit}
 f_n\longrightarrow h\quad\text{locally uniformly on }G,
 \qquad h(G)\text{ is a circle domain}.
\end{equation}
All omitted disks of the finite domains lie in one fixed disk $D_L$.
For every fixed source dust contour, its bounded target sides also
lie in a fixed disk.  The omitted-set bound depends only on a fixed
exterior radius for $G$; the contour bound may also depend on that
fixed contour.  Neither depends on the number of retained tests.

Here the quotient sphere in the scalar corollary is the source sphere
itself, because every source complementary component is a point.
Thus the prescribed $u_i$ are continuous on that quotient, smooth on $G$,
and of finite Dirichlet energy there.  Their boundary value at $b$ is
the single number $u_i(b)$, and condition~\eqref{eq:pair-null} is
precisely the corollary's pairwise trace hypothesis.  The corollary
retains all these original tests while adding auxiliary tests for the
circle-domain selection.  Consequently the core agreement, energy and
oscillation bounds, and contour assertions above hold simultaneously
for the prescribed list.  No countability assumption on $B$ is needed.

It remains to turn these finite transferred tests into their prescribed
continuous Sobolev continuations on the limiting sphere.  That passage,
including the values on positive-area point dust, is proved next.

\subsection{Filling holes and identifying the limiting tests}

We record explicitly the small error introduced by filling round
holes.  Suppose $v$ is bounded and locally Sobolev on a finite circle
domain containing infinity, has finite energy, and has cluster
interval of length at most $A(T)$ at each hole $T$.  Fill $T$ by a
constant in its cluster interval, call the resulting function
$\widetilde v$, and extend $\nabla v$ by zero on the holes to a vector
field $g$.  If $r(T)$ is the radius of $T$, then for
$j=1,2$ and $\varphi\in C_c^\infty(\C)$,
\begin{equation}\label{eq:model-jump}
 \left|\int\widetilde v\,\partial_j\varphi
             +\int g_j\varphi\right|
 \le4\|\varphi\|_\infty\sum_T r(T)A(T).
\end{equation}
Indeed, on almost every line parallel to the $j$th axis the function
is locally absolutely continuous on each domain interval and has finite
variation on its bounded portions relevant to the support of the test
function.  At either endpoint of a hole the jump to its
chosen constant is at most $A(T)$.  Integrating the two endpoint
errors over the transverse projection of that hole, of length
$2r(T)$, gives \eqref{eq:model-jump}.  Point holes affect only finitely
many transverse parameters.  Compactly supported test functions
remove any issue at infinity.

Fix one original smooth test $u=u_i$.  Clip $v_{i,n}$ to a fixed
closed interval containing the range of $u$.  This preserves
\eqref{eq:model-core}, decreases energy and cluster oscillations,
and supplies a uniform bound.  Fill the holes as above, obtaining
$\widetilde v_n$, and write $g_n$ for the domain gradient extended by zero over the holes.
Disjointness and the common target bound imply
\[
 \sum_T r_n(T)^2\le L^2,\qquad
 \sum_T r_n(T)A_{i,n}(T)\le L\varepsilon_n.
\]
Moreover, writing $E_u=\int_G|\nabla u|^2$, we have
\begin{equation}\label{eq:model-energy-bound}
 \int|g_n|^2\le E_u+\varepsilon_n^2.
\end{equation}
Take weak subsequential limits of the bounded functions in spherical
$L^2$ and of their differentials in $L^2$.  Equation
\eqref{eq:model-jump} shows that the limiting differential is the weak
derivative of the limiting function $V$.  This argument takes place in
the planar chart containing all holes; near infinity the retained
core gives the ordinary Sobolev derivatives.  Hence
$V\in W^{1,2}(\Chat)$ and $\int|\nabla V|^2\le E_u$.
On compact subsets of $h(G)$, inverse convergence and
\eqref{eq:model-core} give
\[
 V=u\circ h^{-1}\quad\text{a.e. on }h(G).
\]
By conformal invariance the energy of this function on $h(G)$ is
already $E_u$.  Thus
\begin{equation}\label{eq:model-energy-saturation}
 \nabla V=0\quad\text{a.e. on }F:=\Chat\setminus h(G).
\end{equation}

Zero gradient alone does not identify values on a positive-area set
of point components.  We now determine those values.  Fix a source
dust disk $U$ with smooth boundary $C\subset G$, and an interval
$[\alpha,\beta]$ containing $u(\overline U)$.  Retain $C$ and a collar
in every sufficiently late core.  Let $O_n$ and $O$ be the bounded
sides of $f_n(C)$ and $h(C)$, respectively.  Define
\[
 e_n^+=\begin{cases}(\widetilde v_n-\beta)_+&\text{on }O_n,\\
                         0&\text{off }O_n,\end{cases}
 \qquad
 e_n^-=\begin{cases}(\alpha-\widetilde v_n)_+&\text{on }O_n,\\
                         0&\text{off }O_n.\end{cases}
\]
Every retained-core point mapping into $O_n$ comes from $U$, by the
contour-side correspondence.  Both excesses therefore vanish on
those core pieces and on an inner collar of $f_n(C)$.  Extending them
by zero creates no jump across that contour.  Their extended gradient
norms are at most $\varepsilon_n$, by \eqref{eq:model-tail}, and their
hole-jump errors satisfy \eqref{eq:model-jump}, since positive-part
maps are $1$-Lipschitz.  They are uniformly bounded and supported in
one fixed disk.  Every weak limit consequently has zero
distributional gradient and bounded support, and is zero.

To pass this conclusion to $V$, take a nonnegative smooth function
$\varphi$ compactly supported in $O$.  Winding-number stability puts
its support in $O_n$ eventually.  Therefore
\[
 \int\varphi(\widetilde v_n-\beta)\le\int\varphi e_n^+\longrightarrow0,
 \qquad
 \int\varphi(\alpha-\widetilde v_n)\le\int\varphi e_n^-\longrightarrow0.
\]
Weak convergence gives $\alpha\le V\le\beta$ almost everywhere on
$O$.  In particular, we have not interchanged a nonlinear operation
with a weak limit.

For every integer $\ell\ge1$, choose one finite source dust-disk
cover of $B$ so fine that $u$ oscillates by less than $1/\ell$ on each
closed disk.  Apply the preceding conclusion to those finitely many
intervals.  The union of all their exceptional null sets, over all
$\ell$, is still null.  If $w\in F$ is outside this union and its
source end is $b$, the whole component containing $w$ lies inside
the image contour of each cover disk containing $b$.  It follows that
$|V(w)-u(b)|\le1/\ell$ for every $\ell$, and hence $V(w)=u(b)$.
This also identifies the values on a collection of point components
of positive total area.

The continuous function $u\circ\pi_h$ now agrees almost everywhere
with $V$.  It is the required Sobolev representative, with the
prescribed value on \emph{every} complementary component, and
\eqref{eq:model-energy-saturation} gives its zero gradient there.
The argument applies to each of the finitely many tests.  Finally,
the nonpoint components of $F$ are disjoint round disks in a bounded
region.  Their squared radii have finite sum, so they are countable
and form a null family.  This proves Proposition~\ref{prop:models}
for smooth tests.

\subsection{Exact traces under smoothing}

We use the following elementary extension fact.
\begin{lemma}\label{lem:zero-trace-extension}
Let $F\subset\Chat$ be closed.  If $w\in C(\Chat,\R)$ vanishes on $F$,
is locally Sobolev on $\Chat\setminus F$, and has finite Dirichlet
energy there, then $w\in W^{1,2}(\Chat)$ and
$\nabla w=0$ almost everywhere on $F$.
\end{lemma}
\begin{proof}
For $\delta>0$ the signed truncation
$w_\delta=\operatorname{sgn}(w)(|w|-\delta)_+$ has compact support
away from $F$.  It is globally Sobolev by localization, and its
gradient norm is at most that of $w$ off $F$; its gradient is zero
on $F$.  Since $\|w_\delta-w\|_\infty\le\delta$, weak compactness
proves the assertion, including the vanishing gradient on $F$.
\end{proof}

For a continuous test $u$ as in the proposition, choose a locally
finite smooth partition of unity $(\rho_k)$ on $G$, with each support
compact in $G$.  Local mollification gives smooth functions $a_k$
near these supports, with both uniform error and local gradient
error at most $\eta_k$.  Choose the positive $\eta_k$ so that
\[
 \eta_k\longrightarrow0,\qquad
 \sum_k\eta_k\bigl(1+\|\nabla\rho_k\|_2\bigr)<\infty.
\]
Then $\widehat u=\sum_k\rho_k a_k$ is bounded and smooth on $G$,
and the product rule shows that
$\nabla\widehat u-\nabla u\in L^2(G)$.  As a point approaches $B$,
every fixed finite collection of supports is eventually avoided.
The bound $|\widehat u-u|\le\max\{\eta_k:\rho_k\ne0\}$ therefore
tends to zero there.  Defining $\widehat u=u$ on $B$ gives a continuous
function on the sphere with exactly the original trace.  The same
construction in a chart at infinity is included in the partition.

Apply this to each test.  Their traces and the pairwise condition
\eqref{eq:pair-null} are unchanged.
The smooth case provides one circle-domain model $h$ for the
$\widehat u_i$.  Each function
\[
 w_i=(u_i-\widehat u_i)\circ\pi_h
\]
is continuous, vanishes on $F$, is locally Sobolev on $h(G)$, and has
finite energy there by conformal invariance.  Lemma
\ref{lem:zero-trace-extension} makes it globally Sobolev with zero
gradient on $F$.  Adding it to $\widehat u_i\circ\pi_h$ proves the
model properties for the original $u_i$.  If $B$ is empty, the
identity is already a model.  Proposition~\ref{prop:models} follows
in all cases.
